% Français 9115, batch 25 — Gallica views 481–500.
% Pass: Opus 5.5 (claude-opus-5-5), 2026-09-22 — first pass, unchecked against the views by a human.
%
% What the twenty views hold. One piece, begun before the batch and going
% on after it, in French, in her regular fair hand, fairly heavily
% corrected (words struck on the line and rewritten above it, a few
% marginal additions): the draft of her prize memoir on vibrating elastic
% surfaces, in its section comparing theory with Chladni's experiments on
% square plates — first the nodal figures, then (N.o 16, v. 491) the
% sounds. Pages (73) to (90) in her pagination, written on both sides of
% each leaf; one slip of calculations bound between (74) and (75). No
% microfilm target, no printed matter, no number theory.
%
%   v. 481  (73) the nodal lines of Chladni's fig. 70 agree with the
%           theory; the sounds of figs 69 and 70 proportional to
%           −m²(Π″ − 2) = 10; fig. 87 b (m = 2, b = −B): four lines from
%           the periodic factor, six diagonal lines, four forming a
%           perfect square; figs 87 a and 88 a: four sinuous lines with
%           two complete flexions each.
%   v. 482  (74) fig. 88 b: cos²(π.2x/A) = (b+B)/4b, cos²(π.2y/A) =
%           (b+B)/4B, sixteen lines from the periodic factor; the guess
%           that the figure belongs to b = B.
%   v. 483  a slip of calculations without a word (ink 237): the
%           variation of a double integral of ds divided by
%           (1/R + 1/R′)², with the curvature written out in dz/dx,
%           dz/dy, ddz/dx², ddz/dxdy, ddz/dy², and quantities N, P, P′,
%           Q, Q′, Q″. It is not a page of the memoir.
%   v. 484  the blank back of the slip's mounting sheet, a few blots and
%           the writing of v. 483 showing through: not transcribed.
%   v. 485  (75) end of the sentence of v. 482; sounds of figs 87 b, 87 a,
%           88 a, 88 b proportional to 40; with Π″ = 1 − 25 instead of
%           1 − 9, integral (G) becomes a new z; for m = 1, b = −B, two
%           diagonal lines and four curved ones ending where cos = ±√3/2;
%           fig. 79 a, sound proportional to 26.
%   v. 486  (76) she leaves integral (G); praise of Chladni for having
%           guessed which distortions of his figures were due to
%           imperfections of the plate.
%   v. 487  (77) integral (H): straight nodal lines, sounds following the
%           law of the lamina supported at both ends or free with dz/dx =
%           0, d³z/dx³ = 0; its series never terminate; integral (I),
%           periodic in y only, the same singularity.
%   v. 488  (78) figs 79 b and 86 referred to (I) with m = 5, m = 6; the
%           case m = 5 worked: a line parallel to the axis of y, two fixed
%           points of rest, two flexions; fig. 79 b fits except for three
%           lines she takes for equivalent.
%   v. 489  (79) the case m = 6 and fig. 86, the two re-entrant curves
%           taken for an accident of the plate.
%   v. 490  80 her explanation of 79 b and 86 « conjecturale »; sounds 25
%           and 36; she has accounted for more than half of Chladni's
%           square-plate figures; she will follow Chladni's own
%           equivalences for the sounds of the rest.
%   v. 491  (81) N.o 16 « comparaison des sons »: the effect of a small
%           inequality of thickness; Chladni's own bound on experimental
%           error, quoted: never more than a semitone (« p 154 de son
%           livre »), hence at most a tone between two sounds.
%   v. 492  (82) the « vibrations tournantes » (half-periods only): the
%           theoretical numbers for figs 63, 66 a, 69–70, 74, 79 a, 79 b,
%           86 (2, 5, 10, 17, 26, 25, 36), their ratios to fig. 63, and
%           the sounds of « la table p. 152 du traité d'acoustique »; 63
%           to 66 a: an octave, a major and a minor tone, = 5/2.
%   v. 493  (83) 63 to 69 agrees (5); fig. 70 a tone off fig. 69, which
%           she cannot explain; 63 to 74 within the error; 63 to 79 a,
%           « beaucoup trop petit ».
%   v. 494  (84) 79 a about a tone too sharp, 79 b worse; 63 to 86, 64/3
%           against 18; her explanations of 79 b and 86 only
%           « hypothetique »; she turns to the figures with complete
%           periods in both directions.
%   v. 495  (85) the same comparison against fig. 68 a (8): figs 71, 75,
%           80 a, 81 a, 82, 87–88, 79 a (13, 18, 25, 29, 32, 40, 26),
%           ratios, and Chladni's table.
%   v. 496  86 each ratio checked: 71 and 75 agree, 80 differs by a
%           comma, 81 by more than a comma, 82 and 87–88 agree.
%   v. 497  (87) 68 to 79 agrees within the comma 128/125; the cases for
%           which Chladni gave no figure, in his notation m/n (m straight
%           nodal lines parallel to one axis, n to the other); theoretical
%           numbers for 2/2, 4/2, 5/3, 6/3, 5/4, 6/4, 5/5, 6/5.
%   v. 498  (88) their ratios to 2/2 and Chladni's sounds; 4/2 and 5/3
%           agree taking the reading most favourable to the theory.
%   v. 499  (89) 6/3, 5/4, 6/4, 5/5, 6/5 each within a semitone or a
%           comma.
%   v. 500  (90) « Il étoit impossible d'espérer plus d'accord entre la
%           théorie et l'experience »; Chladni himself took as equivalent
%           figures whose sounds differ by a semitone, a tone or a third
%           (« V. latable p. 152 »); « Cependant il se présente ici une
%           difficulté: »… The page stops mid-sentence on « dans le sens
%           des x ».
%
% Seams. View 481 opens a new paragraph on page (73), « On voit donc que
% la théorie est parfaitement d'accord… » (its thumbnail was compared with
% that of v. 480, page (72), which ends on « et d'autre des axes. »), so
% nothing is cut at 480/481 — but the section on the square plates began
% before this batch, and figs 69, 70, 87 b and the integrals (C)–(I) are
% introduced there. At 500/501 the sentence is cut: v. 501, page (91),
% ink 246, opens « et dans celui des y, lorsque on prend pour terme de
% comparaison… ». References to pages outside the batch: N.o 13 and
% N.o 15 (v. 492, 493, 495), « comme je l'ai expliqué à l'occasion de la
% fig 70 » (v. 481), « comme je l'ai dit plus haut » (v. 485), the
% integrals (G), (H), (I).
%
% The numbers on the leaves, and why no \folio{} is written. (1) Her
% pagination at the head, centre: (73) v. 481, (74) v. 482, (75)–(79)
% v. 485–489, 80 v. 490, (81)–(85) v. 491–495, 86 v. 496, (87)–(90)
% v. 497–500 — without a gap; 80 and 86 are written without parentheses;
% (73), (75) and (77) are struck through with a stroke, the others not; in
% (89) the 8 is written over another figure. The slip of v. 483 has no
% page number. (2) The ink series at the top right, in the larger rounder
% hand: 236 (v. 481), 237 (v. 483, the slip), 238 (v. 485), 239 (v. 487),
% 240 (v. 489), 241 (v. 491), 242 (v. 493), 243 (v. 495), 244 (v. 497),
% 245 (v. 499) — one per leaf, on the side her pagination makes odd, never
% on the even-paginated back, the slip counted as a leaf; 246 at v. 501
% continues it. As in every batch of this volume it is penned, not
% pencilled, and no \folio{} is written. (3) Her paragraph numbers:
% N.o 16 (v. 491), and the references N.o 13, N.o 15. (4) Chladni's figure
% numbers (63–88, with a and b) and his page « p. 152 », « p 154 ». Every
% number above was read on its view; none is inferred. None of these
% leaves can be Laubenbacher and Pengelley's « Manuscript C »
% (ff. 348r–349r): the subject is the elastic plate, and the ink series
% here runs 236–245.
%
% Notation. Hers, kept. A the side of the square plate; m, n the numbers
% of periods; b, B, a′ and a looped cursive capital distinct from A, set
% $\mathcal{A}$ as in batch 22; the tall-legged letter of −m²(Π″ − 2) is
% set Π″, distinct from the π of the cosines, as batch 22 did. Her
% arguments of sine and cosine are written in brackets over A, « cos²
% (π.2x)/A », and are set \frac{(\pi.2x)}{A}; « sin » and « cos » are set
% \text{sin}, \text{cos}. « = … » and « z = … » stand for a factor she
% leaves blank. Musical pitches are Chladni's octave-numbered names
% (« sol.1 », « si♭ 3 », « ut♯ 6 »), set with \flat and \sharp; her tables
% are set as arrays, the leader dots and dashes shortened to \dots.
% Chladni's notation m/n for the square-plate modes is kept as a slash.
% Words underlined in the prose (« comparaison des sons », « comma ») are
% set in \emph{}; letters underlined in mathematics keep \underline{}.
% Where the leaf and the calculation disagree the leaf stands. Spelling is
% hers and is not flagged: « complettes », « complettér », « quarré »,
% « parallelle(s) », « equivalantes », « proportionels », « d'acord »,
% « apperçu », « bizares », « distortions », « emploirai », « précèdantes »,
% « suivans », « précédans », « differens », « mouvemens », « longtems »,
% « paroissent », « seroit », « comma », « demiton », « la figures »,
% « une autre figure nodales », and the run-together words « dela »,
% « delaplaque », « delafigure », « àlaquelle », « lamême », « lamoitié »,
% « ceque », « cequi », « desorte », « adonné », « aulieu », « aucontraire »,
% « audessous », « parconséquent », « ladifference », « latable »,
% « ellemême », « peuplus ».
%
% What could not be read, and what is offered with doubt. \ill{} stands
% eleven times: the blotted sign in −m²(Π″ \ill{} 2) (v. 481); on the
% slip, two heavily struck lines, the first sign of the long fraction and
% a blotted sign inside it (v. 483, four in all); the struck words under
% « deviné » and after « s'est bien gardé de » (v. 486); a blotted sign in
% the series factor (v. 488); the struck numerator on v. 489; a struck
% capital before « pour n'être pas » (v. 491); a struck sign before
% « majeur » (v. 494). \uncertain{} stands for Π in her marginal « je
% changerai le signe des Π » (v. 481); « senti », « la presentoit »
% (v. 486, struck); « Donnés » and « Avant » (v. 488, struck) and b′
% (v. 488); « deux » (v. 492, struck); n′ = 1 (v. 489); « p » in « p 154 »
% (v. 491); N.o 15, « trop aigu mais » and « pourtant » (v. 493); 128/125
% (v. 499); « -lles » (v. 500, struck). Nothing was guessed to fill a gap.
%
% Nothing here was checked against any published edition of Chladni or of
% Sophie Germain's papers.
\documentclass[11pt,a4paper]{article}
% The preamble shared by every transcription of the mathematicians' manuscripts in Gallica.
%
% It rests on one idea: **uncertainty compounds, it does not resolve.** A
% transcription that smooths over a doubtful word has destroyed the only thing
% it was made for — a reader must be able to tell, at every point, what was on
% the page and what was supplied. The apparatus macros below are therefore the
% critical apparatus, not decoration.
%
% The same commands are recognised by `scripts/render.mjs`, which produces the
% site's reading view, and by `scripts/tei.mjs`. A macro added here must be
% added there too, or rendering fails loudly — which is the wanted behaviour.
%
% Comments here are English, like the rest of the repository. The documents
% this preamble sets are French, and that is deliberate: see the skills.

\ProvidesPackage{germain}[2026/09/15 Transcriptions of mathematicians' manuscripts in Gallica]

\usepackage{amsmath,amssymb,amsthm}
\usepackage{mathrsfs}
% Kept from the parent preamble so the renderer's diagram support stays
% available; these manuscripts draw no commutative diagrams, and nothing here uses it.
\usepackage{tikz-cd}
\usepackage[english,french]{babel}
\usepackage{fontspec}

% --- Greek ------------------------------------------------------------------
% The text face stays fontspec's default, Latin Modern, which has no Greek:
% XeTeX drops every glyph it lacks with a line in the log and nothing on the
% page, and the Greek words the manuscripts quote vanished from the PDFs. So
% the Greek and Greek Extended blocks get a XeTeX character class of their own,
% and entering or leaving a run of it switches the family to CMU Serif — the
% same Computer Modern design, with polytonic Greek — and back. Only the
% family changes: a Greek word in \textit or \textbf keeps its shape and series.
% The transitions to and from every other class carry the tokens that class
% has towards an ordinary letter, so French spacing before « ; » or inside
% « » still holds after a Greek word. No grouping, so no brace or \endgroup
% can come out unbalanced; maths is left alone, since XeTeX inserts nothing
% there. Kept identical in `exercices.sty`.
\makeatletter
\newfontfamily\germainGreekfont{cmun}[
  Extension=.otf, NFSSFamily=germaingreek,
  UprightFont=*rm, ItalicFont=*ti, SlantedFont=*sl,
  BoldFont=*bx, BoldItalicFont=*bi, BoldSlantedFont=*bl]
\newXeTeXintercharclass\germain@greek
\def\germain@greekrange#1#2{%
  \@tempcnta=#1\relax
  \loop
    \XeTeXcharclass\@tempcnta=\germain@greek
    \ifnum\@tempcnta<#2\advance\@tempcnta\@ne\repeat}
\germain@greekrange{"0370}{"03FF}
\germain@greekrange{"1F00}{"1FFF}
\def\germain@greekprev{\rmdefault}
\def\germain@greekin{%
  \edef\germain@greekprev{\f@family}\fontfamily{germaingreek}\selectfont}
\def\germain@greekout{\fontfamily{\germain@greekprev}\selectfont}
\def\germain@greekjoin{%
  \XeTeXinterchartokenstate=\@ne
  \@tempcnta=\z@
  \loop
    \ifnum\@tempcnta=\germain@greek\else
      \XeTeXinterchartoks\germain@greek\@tempcnta=\expandafter{%
        \expandafter\germain@greekout\the\XeTeXinterchartoks\z@\@tempcnta}%
      \XeTeXinterchartoks\@tempcnta\germain@greek=\expandafter{%
        \the\XeTeXinterchartoks\@tempcnta\z@\germain@greekin}%
    \fi
    \ifnum\@tempcnta<4095\advance\@tempcnta\@ne\repeat}
% babel-french sets its own classes, and saves and restores the interchar
% state, each time a language is selected: join again after it does.
\addto\extrasfrench{\germain@greekjoin}
\addto\extrasenglish{\germain@greekjoin}
\AtBeginDocument{\germain@greekjoin}
\makeatother

\usepackage{geometry}
\geometry{a4paper,margin=25mm,marginparwidth=28mm}
\usepackage{marginnote}
\usepackage{xcolor}
\usepackage[normalem]{ulem}
\setlength{\emergencystretch}{3em}
\usepackage{hyperref}
\hypersetup{colorlinks=true,linkcolor=black,urlcolor=black,pdfborder={0 0 0}}

% --- Batch metadata ---------------------------------------------------------
% Taken as they stand by the rendering script, which shows them at the head of
% the reading view. They fix what the file claims to cover. The macro names are
% the parent project's, so that the scripts stay shared; \folder here means the
% volume's slug — `fr-9115` — and \pages counts Gallica views.

\newcommand{\folder}[1]{\gdef\grothFolder{#1}}
\newcommand{\batch}[1]{\gdef\grothBatch{#1}}
\newcommand{\pages}[2]{\gdef\grothFirst{#1}\gdef\grothLast{#2}}
\newcommand{\foldertitle}[1]{\gdef\grothFoldertitle{#1}}

% The holder's own shelfmark, printed as they write it: « Français 9115 ».
\newcommand{\shelfmark}[1]{\gdef\grothShelfmark{#1}}

% The dating, copied verbatim from the holder's catalogue — never inferred
% here. « XIXe siècle » is the BnF's; a pass that dates a leaf from its content
% says so in a \note{}, as its own inference.
\newcommand{\dating}[1]{\gdef\grothDating{#1}}

% The status notice, printed in the title block and repeated as a diagonal
% watermark on every page of the PDF. The manuscripts are in the public
% domain; what this line declares is not a right but a quality — an
% unverified first pass by a machine. The skills write the line; a file
% without it gets no watermark, so leaving it out is visible in the output.
\newcommand{\watermark}[1]{\gdef\grothWatermark{#1}}

\gdef\grothFolder{?}\gdef\grothBatch{}\gdef\grothFirst{?}\gdef\grothLast{?}%
\gdef\grothFoldertitle{}\gdef\grothShelfmark{}\gdef\grothDating{}\gdef\grothWatermark{}

% --- The résumé -------------------------------------------------------------
\newenvironment{resume}
  {\small\setlength{\parskip}{0.45em}}
  {\par\medskip\hrule\medskip}

% --- Keywords ---------------------------------------------------------------
% In English, closing the modernised reading's résumé: the modern vocabulary
% under which to search for what the leaves construct. The manifest extracts
% them to tag the volume — this line is the single source of the tags.
\newcommand{\keywords}[1]{\par\smallskip\noindent
  {\footnotesize\textsc{Keywords} --- \textit{#1}}\par}

% --- The critical apparatus -------------------------------------------------

% The Gallica view number, in the margin. This is the anchor that lets a
% disagreement over a formula be settled by looking at *one* image rather than
% a volume of 750. The site uses it to turn the facsimile as the transcription
% is scrolled. A view is one scanned image: usually one side of a leaf.
\newcommand{\page}[1]{%
  \marginnote{\footnotesize\color{gray!70}\textsf{v.\,#1}}%
  \label{page:#1}\ignorespaces}

% The BnF's pencilled foliation, where the leaf carries it and a pass can read
% it: « 198r ». This is what the literature cites — « ff. 198r–208v » — and
% the one line that turns a citation into a view. Recorded, never inferred:
% a folio a pass cannot read is not written.
\newcommand{\folio}[1]{%
  \marginnote{\footnotesize\color{gray!50}\textsf{f.\,#1}}\ignorespaces}

% The modernised reading's coarser counterpart to \page{}: which views a
% section is drawn from.
\newcommand{\pagerange}[2]{%
  \marginnote{\footnotesize\color{gray!70}\textsf{v.\,#1--#2}}%
  \ignorespaces}

% Illegible: never guessed.
\newcommand{\ill}{\textcolor{red!60!black}{[\,\ldots\,]}}

% A doubtful reading: the word is given, and flagged as uncertain.
\newcommand{\uncertain}[1]{\uline{#1}}

% An editorial addition — a missing letter, an implied word.
\newcommand{\add}[1]{\textcolor{blue!50!black}{[#1]}}

% Struck out by the writer. What was crossed out is part of the document.
\newcommand{\struck}[1]{\sout{#1}}

% The transcriber's note, on the state of the page or the mathematics.
\newcommand{\note}[1]{\par\smallskip{\small\color{gray!60!black}\textit{[#1]}}\par\smallskip}

% A marginal note of hers, distinct from the body of the text.
\newcommand{\marginal}[1]{\par\smallskip\noindent\hspace{1em}%
  {\small\color{gray!60!black}#1}\par\smallskip}

% --- Title ------------------------------------------------------------------
% The title block is French, like the documents themselves.

\AddToHook{shipout/background}{%
  \ifx\grothWatermark\empty\else
    \begin{tikzpicture}[remember picture, overlay]
      \node[rotate=52, scale=3.4, align=center, text=gray!18,
            font=\sffamily\bfseries]
        at (current page.center) {\grothWatermark};
    \end{tikzpicture}%
  \fi}

\AtBeginDocument{%
  \begin{center}
    {\large\bfseries Manuscrits de mathématiciens (Gallica) ---
      \ifx\grothShelfmark\empty\grothFolder\else\grothShelfmark\fi}\\[2pt]
    {\itshape\grothFoldertitle}\\[4pt]
    {\small vues \grothFirst--\grothLast
      \ifx\grothBatch\empty\else\ (lot \grothBatch)\fi}\\[2pt]
    \ifx\grothDating\empty\else
      {\small\color{gray!60!black}Datation du catalogue : \grothDating}\\[2pt]
    \fi
    \ifx\grothWatermark\empty\else
      {\small\bfseries\color{red!45!gray}\grothWatermark}\\[2pt]
    \fi
    {\footnotesize\color{gray!60!black}Transcription établie d'après les images IIIF de Gallica.
     Source gallica.bnf.fr / Bibliothèque nationale de France.\\
     Les lectures incertaines sont soulignées, les illisibles notées [\,\ldots\,],
     les ajouts entre crochets ; \textsf{v.} numérote les vues de Gallica,
     \textsf{f.} la foliotation au crayon de la BnF.}
  \end{center}
  \vspace{1em}}

\endinput


\folder{fr-9115}
\shelfmark{Français 9115}
\batch{25}
\pages{481}{500}
\dating{1801-1900}
\watermark{Lecture automatique\\non vérifiée}
\foldertitle{Papiers de Sophie GERMAIN. — Recueil de dissertations et problèmes mathématiques et physiques.}

\begin{document}

\note{Numérotation des feuillets. Chaque page de ce brouillon porte en
tête, au milieu, sa pagination : (73) et (74) aux vues 481 et 482, (75) à
(79) aux vues 485 à 489, 80 à la vue 490, (81) à (85) aux vues 491 à 495,
86 à la vue 496, (87) à (90) aux vues 497 à 500, sans lacune ; 80 et 86
sont écrits sans parenthèses ; (73), (75) et (77) sont barrés d'un trait.
Les feuillets sont écrits des deux côtés, et seul le côté de pagination
impaire porte en haut à droite, à l'encre, d'une écriture plus ronde, un
nombre de la série qui court dans tout le volume : 236 (vue 481), 238
(485), 239 (487), 240 (489), 241 (491), 242 (493), 243 (495), 244 (497),
245 (499) ; le petit feuillet de calculs de la vue 483 porte 237. Cette
série numérote des feuillets et non des pages, et elle compte le petit
feuillet ; elle occupe la place de la foliotation de la Bibliothèque,
mais elle est tracée à la plume et non au crayon, et aucun « f. » n'est
donc porté. Tous ces nombres ont été lus sur la vue ; aucun n'est déduit.
La vue 484 est le dos blanc de la feuille qui porte le feuillet de la vue
483 ; elle n'est pas transcrite.}

\page{481}

On voit donc que la théorie est parfaitement d'accord avec l'experi-
ence sur la disposition des lignes nodales qui composent
la fig 70.

Les tons \add{qui accompagnent} \struck{des} fig. 69 et 70 doivent être proportionnels
au nombre $-m^{2}(\Pi''\,\ill{}\,2) = 10$. \add{+ dorenavant je changerai le signe des \uncertain{$\Pi$}.}

Pour la fig. 87 b $m = 2$ $b = -B$. Conformement à l'experience, il doit y avoir
\struck{outre les} quatre lignes nodales dépendantes du facteur periodique
\struck{lesquelles forment entre elles quatre intersections qui sont comme
on sait autant de points d'origine}; et parceque les quatre
intersections formés par ces lignes, sont, comme on sait, autant de
points d'origine; le facteur $\text{cos}^{2}\frac{(\pi.2x)}{A} - \text{cos}^{2}\frac{(\pi.2y)}{A}$ doit donner, comme
il donne en effet, deux lignes nodales diagonales passant par
chacun de ces points. Parmi les lignes nodales diagonales, il
y en a deux qui passent chacune par deux des points
d'origines; de sorte que leur nombre se réduit a six; les
quatre autres forment un quarré parfait comme le montre
la figure donnée par M\textsuperscript{r} Chladni

Pour les fig. 87 a et 88 a $m = 2$. \struck{mais la relation entre $b$ et $B$
donne lieu}, comme je l'ai expliqué à l'occasion de la fig 70,
la relation entre $b$ et $B$ donne lieu à la formation de lignes
sinueuses. Ces fig. renferment deux périodes complettes dans le sens
des deux axes: ainsi il doit y avoir, comme il y a en effet,
quatre lignes sinueuses; et chacune de ces lignes devant passer
par quatre points invariables de repos, \struck{elle} doit \struck{montrer}
eprouver, \struck{comme l'experience l'a montré}, deux flexions complettes.
c'est aussi ceque l'experience adonné comme le montrent
les fig dont il s'agit.
\note{En tête, au milieu, « (73) » barré d'un trait ; en haut à droite, à
l'encre, d'une écriture plus ronde, « 236 ». La lettre à hautes jambes suivie de deux traits, dans
« $-m^{2}(\Pi''\,\ill{}\,2)$ », est la même qu'aux vues 485 et suivantes
(« $\Pi'' = 1 - 25$ au lieu de $\Pi'' = 1 - 9$ ») ; on la note $\Pi''$ comme
les lots précédents, pour la distinguer du $\pi$ des cosinus. Le signe
entre $\Pi''$ et le 2 est noyé sous une tache d'encre ; le moins qui précède $m^{2}$ est lié au mot « nombre » par un trait de
plume, comme à la vue 485 ; le signe d'égalité
qui suit est repassé en un trait épais. La remarque « + dorenavant je
changerai le signe des \uncertain{$\Pi$}. », à la suite de « 10. » sur la
même ligne, est d'une plume plus fine ; la croix qui l'ouvre est la
sienne ; après « des » vient un signe à deux jambages, semblable à son
$\Pi$ mais qu'on pourrait lire « 17 », suivi d'un tiret. Les deux passages donnés pour barrés
(« outre les », « lesquelles forment … d'origine ») sont traversés par des
traits qui coupent les mots à mi-hauteur, distincts des traits de liaison
de sa main ; le second est repris aussitôt, presque mot pour mot, par
« et parceque les quatre intersections … ». Dans les deux cosinus, le 2
devant $x$ est écrit sur un autre chiffre ; le $\pi$ est son $\pi$ à hautes
jambes. Le bas de la page est blanc.}

\page{482}

Pour la fig. 88 b. $m = 2$. La relation entre $b$ et $B$ permet de remplir
les deux conditions $\text{cos}^{2}\frac{(\pi.2x)}{A} = \frac{b+B}{4b}$ et $\text{cos}^{2}\frac{(\pi.2y)}{A} = \frac{b+B}{4B}$.
ainsi il doit y avoir, comme il y a en effet, \struck{4 lignes
nodales dependantes du fa.} autour de chaque point d'origine
4 lignes nodales dependantes du facteur $\text{cos}^{2}\frac{(\pi.2x)}{A} - \text{cos}^{2}\frac{(\pi.2y)}{A} = 0$;
et parceque cette fig. renferme comme les précédentes 2 periodes
complettes prise dans le sens des deux axes, ce qui fait en tout
4 periodes puisque chacune est prise deux fois, il doit y avoir
16 lignes dues à ce facteur. Il est clair que quatre d'entre elles
doivent former une courbe rentrante autour du centre dela
figure, et que deux seulement doivent finir dans les \struck{periodes
complettes dans le sens d'un des axes et moitié de periodes dans
l'autre sens} parties delafigure qui sont a 2 limites et
ne sont composées parconséquent que d'une demie periode
dans le sens d'un axe, quoiqu'elles renferment une periode
complette dans le sens de l'autre axe.

Tout ceque je viens de dire, est indépendant des valeurs particul\add{ières}
de $b$ et $B$, et se trouve parfaitement d'acord avec l'experience
dans le cas de la fig 88 b. je serois même tenté de croire
que cette fig est due à la relation $b = B$: car alors
$\text{cos}\frac{(\pi.2x)}{A}$ et $\text{cos}\frac{(\pi.2y)}{A}$ seroient egaux entre eux, et à $\sqrt{\tfrac{1}{2}}$: cequi
donneroit des lignes nodales \struck{semblement} \add{semblablement} situées par rapport
à l'un et l'autre des axes, d'une courbure uniforme
dans leur deux moitiés, et dont l'origine seroit à lamoitié
\note{En tête, au milieu, « (74) », non barré, le 4 repassé ; aucun nombre
en haut à droite. Au-dessous de « effet, » un petit signe, peut-être un
renvoi, sans note qui lui réponde. Le « 4 » qui ouvre la cinquième ligne
est tracé d'un seul trait qui le traverse, comme son 4 ailleurs. Dans
« delafigure », une petite croix sur la finale annule un pluriel ; dans
« a 2 limites », le 2 est suivi d'une boucle qui peut être un s. Le mot
« particul\add{ières} » est coupé au bord du feuillet, la fin passant sous
le montage. « semblement » est barré et « semblablement » écrit au-dessus.
La phrase s'arrête au bas de la page sur « à lamoitié » ; elle reprend à la
vue 485, page (75), après le feuillet de calculs de la vue 483 et la vue
484.}

\page{483}

\[ \frac{\struck{1+d.}}{\left(\frac{1}{R}+\frac{1}{R'}\right)}^{2} = -\frac{h}{k^{3}} \]
\struck{\ill{}}

\struck{\ill{}}
\[ = \frac{-\left(1+\left(\frac{dz}{dy}\right)^{2}\right)\frac{ddz}{dx^{2}} + 2.\frac{dz}{dx}.\frac{dz}{dy}.\frac{ddz}{dxdy} - \left(1+\left(\frac{dz}{dx}\right)^{2}\right)\frac{ddz}{dy^{2}}}{\sqrt{(1+p^{2}+q^{2})^{3}}} \]
\[ \frac{ds}{\left(\frac{1}{R}+\frac{1}{R'}\right)^{2}} = -\frac{\left\{\left(1+\left(\frac{dz}{dy}\right)^{2}\right)\frac{ddz}{dx^{2}} - 2.\frac{dz}{dx}.\frac{dz}{dy}.\frac{ddz}{dxdy} + \left(1+\left(\frac{dz}{dx}\right)^{2}\right)\frac{ddz}{dy^{2}}\right\}^{2}dxdy}{1+\left(\frac{dz}{dy}\right)^{2}+\left(\frac{dz}{dx}\right)^{2}} \]
\[ \iint\delta.\frac{ds}{\left(\frac{1}{R}+\frac{1}{R'}\right)^{2}} = 2\frac{\left\{\left(1+\left(\frac{dz}{dy}\right)^{2}\right)\frac{ddz}{dx^{2}} - 2.\frac{dz}{dx}.\frac{dz}{dy}.\frac{ddz}{dxdy} + \left(1+\left(\frac{dz}{dx}\right)^{2}\right)\frac{ddz}{dy^{2}}\right\}}{1+\left(\frac{dz}{dy}\right)^{2}+\left(\frac{dz}{dx}\right)^{2}} \]
\[ \frac{\begin{array}{l} \ill{}\left(1+\left(\frac{dz}{dy}\right)^{2}+\left(\frac{dz}{dx}\right)^{2}\right)\left\{2.\frac{dz}{dy}.\frac{\delta dz}{dy}.\frac{ddz}{dx^{2}} - 2\frac{dz}{dx}.\frac{\delta.dz}{dy}.\frac{ddz}{dxdy}\right. \\ \quad \left. + 2.\frac{dz}{dx}.\frac{\delta.dz}{dx}.\frac{ddz}{dy^{2}} - 2.\frac{dz}{dy}.\frac{\delta.dz}{dx}.\frac{ddz}{dxdy}\right\} \\ \quad \left(1+\left(\frac{dz}{dy}\right)^{2}\right)\frac{dd\delta z}{dx^{2}} - 2.\frac{dz}{dx}.\frac{dz}{dy}.\frac{\delta ddz}{dxdy} \ill{} \left(1+\left(\frac{dz}{dx}\right)^{2}\right)\frac{\delta.ddz}{dy^{2}} \end{array}}{\left\{1+\left(\frac{dz}{dy}\right)^{2}+\left(\frac{dz}{dx}\right)^{2}\right\}^{2}} \]
\[ -\frac{\left\{\left(1+\left(\frac{dz}{dy}\right)^{2}\right)\frac{ddz}{dx^{2}} - 2.\frac{dz}{dx}.\frac{dz}{dy}.\frac{ddz}{dxdy} + \left(1+\left(\frac{dz}{dx}\right)^{2}\right)\frac{ddz}{dy^{2}}\right\}\left\{2.\frac{dz}{dy}.\frac{\delta dz}{dy} + 2.\frac{dz}{dx}.\frac{\delta dz}{dx}\right\}}{\left\{1+\left(\frac{dz}{dy}\right)^{2}+\left(\frac{dz}{dx}\right)^{2}\right\}^{2}} \]
\[ N = 0 \qquad P = 2\frac{\left\{\left(1+\left(\frac{dz}{dy}\right)^{2}\right)\frac{ddz}{dx^{2}} - 2.\frac{dz}{dx}.\frac{dz}{dy}.\frac{ddz}{dxdy} + \left(1+\left(\frac{dz}{dx}\right)^{2}\right)\frac{ddz}{dy^{2}}\right\}}{\left\{1+\left(\frac{dz}{dy}\right)^{2}+\left(\frac{dz}{dx}\right)^{2}\right\}^{3}} \]
\[ \begin{array}{l} \underline{\left\{1+\left(\frac{dz}{dy}\right)^{2}+\left(\frac{dz}{dx}\right)^{2}\right\}}\left\{\underline{2\frac{dz}{dx}.\frac{ddz}{dy^{2}}} - 2\frac{dz}{dy}.\frac{ddz}{dxdy}\right\} - \underline{\left(1+\left(\frac{dz}{dx}\right)^{2}\right).2.\frac{dz}{dx}.\frac{ddz}{dy^{2}}} \\ - 2\left(\frac{dz}{dx}\right)\left(1+\left(\frac{dz}{dy}\right)^{2}\right).\frac{ddz}{dx^{2}} + 4.\left(\frac{dz}{dx}\right)^{2}.\frac{dz}{dy}.\frac{ddz}{dxdy} \end{array} \]
\[ = \dots 2\left\{\left(\frac{dz}{dy}\right)^{2}.\frac{dz}{dx}.\frac{ddz}{dy^{2}} - \left(1+\left(\frac{dz}{dy}\right)^{2}\right)\left(\frac{dz}{dx}\right)\frac{ddz}{dx^{2}} - \left(1+\left(\frac{dz}{dy}\right)^{2}-\left(\frac{dz}{dx}\right)^{2}\right)\frac{ddz}{dxdy}\right\} \]
\[ P' = \dots 2\left\{\left(\frac{dz}{dx}\right)^{2}.\frac{dz}{dy}.\frac{ddz}{dx^{2}} - \left(1+\left(\frac{dz}{dx}\right)^{2}\right).\left(\frac{dz}{dy}\right).\frac{ddz}{dy^{2}} - \left(1+\left(\frac{dz}{dx}\right)^{2}-\left(\frac{dz}{dy}\right)^{2}\right)\frac{ddz}{dxdy}\right\} \]
\[ Q = \dots \left(1+\left(\frac{dz}{dy}\right)^{2}\right)\frac{ddz}{dx^{2}}.\left(1+\left(\frac{dz}{dy}\right)^{2}+\left(\frac{dz}{dx}\right)^{2}\right) \]
\[ Q' = \dots -2.\frac{dz}{dx}.\frac{dz}{dy}.\frac{ddz}{dxdy}\left(1+\left(\frac{dz}{dy}\right)^{2}+\left(\frac{dz}{dx}\right)^{2}\right) \]
\[ Q'' = \dots \left(1+\left(\frac{dz}{dx}\right)^{2}\right)\left(1+\left(\frac{dz}{dy}\right)^{2}+\left(\frac{dz}{dx}\right)^{2}\right)\frac{ddz}{dy^{2}} \]
\note{Feuillet de calculs sans un mot, plus petit que les pages du mémoire,
monté sur une feuille ; en haut à droite, à l'encre, « 237 » ; aucune
pagination. Les calculs portent sur la variation d'une intégrale double de
$ds$ divisé par le carré de la somme des courbures
$\frac{1}{R}+\frac{1}{R'}$. À la première ligne, « 1+d. » est écrit et barré
au-dessus du trait qui couvre $\left(\frac{1}{R}+\frac{1}{R'}\right)$, et
l'exposant 2 est placé au-dessus de la parenthèse ; aux lignes suivantes,
$ds$ est écrit de même au-dessus d'un trait : on le donne comme une
fraction, sans décider si le trait est une barre de fraction ou un trait
de réunion. Les deux lignes qui suivent la première (« $= -$ … $+ q^{2}$ … »,
« $=$ … $+$ … ») sont biffées lourdement et illisibles. Dans la grande
fraction, le premier signe est une tache d'encre, et le signe entre
$\frac{\delta ddz}{dxdy}$ et $\left(1+\left(\frac{dz}{dx}\right)^{2}\right)$
est une autre tache ; dans $\frac{dd\delta z}{dx^{2}}$ le $\delta$ est écrit
sur une autre lettre, entre les deux $d$. Les soulignements de la ligne
qui suit $P$ sont les siens. Dans « $+ 4.$ », le 4 est écrit sur un autre
chiffre. Les points de suspension après « $=$ », « $P' =$ », « $Q =$ »,
« $Q' =$ », « $Q'' =$ » sont les siens et tiennent la place d'un facteur
omis. Dans la marge de gauche, des essais de plume : deux « 8 », un
« m », une petite branche feuillue dessinée, et au bas une rangée de
petits « o ». Le bas du feuillet est blanc.}

\page{485}

\struck{dela distance entre le point d'origine} du côté compris entre
la ligne nodale due au facteur periodique et l'extremités àlaquelle
elle est parallele; circonstances qui paroissent d'accord avec
la figure \struck{trouvée} observée par M\textsuperscript{r} Chladni

Les sons qui accompagnent les figures 87 b, 87 a, 88 a
et 88 b, doivent être proportionnels au nombre $-m^{2}(\Pi''-2) = 40$.

Si en conservant les suppositions précèdantes, on \struck{mets} fait
pourtant $\Pi'' = 1 - 25$ aulieu de $\Pi'' = 1 - 9$, l'intégrale
(G) deviendra
\[ z = \dots \text{sin}\frac{(\pi x m)}{A}\text{sin}\frac{(\pi y m)}{A}\left\{b\left\{\text{cos}^{2}\frac{(\pi.mx)}{A} - \frac{16}{3.4}\text{cos}^{4}\frac{(\pi.mx)}{A}\right\}\right. \]
\[ + B\left\{\text{cos}^{2}\frac{(\pi m y)}{A} - \frac{16}{3.4}\text{cos}^{4}\frac{(\pi m y)}{A}\right\} \]
\[ \left. - \frac{b+B}{12}\right\} \]
desorte qu'en prenant $m = 1$ $b = -B$, on auroit
\[ z = \dots \text{sin}\frac{(\pi x m)}{A}\text{sin}\frac{(\pi y m)}{A}\left\{\left(\text{cos}^{2}\frac{(\pi x m)}{A} - \text{cos}^{2}\frac{(\pi y m)}{A}\right)\right. \]
\[ \left.\left(1 - \frac{4}{3}\left(\text{cos}^{2}\frac{(\pi x m)}{A} + \text{cos}^{2}\frac{(\pi y m)}{A}\right)\right)\right\} \]
Il est visible que le facteur $\text{cos}^{2}\frac{(\pi x m)}{A} - \text{cos}^{2}\frac{(\pi y m)}{A}$ doit
\struck{comme pour le} donner, comme je l'ai dit plus haut,
deux lignes nodales diagonales; et l'on voit également
que le facteur $1 - \frac{4}{3}\left(\text{cos}^{2}\frac{(\pi x m)}{A} + \text{cos}^{2}\frac{(\pi y m)}{A}\right)$ devra donner quatre
lignes \struck{courbes} nodales courbes \struck{terminées} à chaque limite
dans les points pour lesquels $\text{cos}\frac{(\pi x m)}{A}$ et $\text{cos}\frac{(\pi y m)}{A}$ seront
$= \pm\sqrt{3}.\frac{1}{2}$.

La fig. 79 a. paroit appartenir à ce cas.
Le son qui l'accompagne, doit donc être
proportionnel au nombre $-m^{2}(\Pi''-2) = 26$.
\note{En tête, au milieu, « (75) » barré d'un trait ; en haut à droite, à
l'encre, « 238 ». Le feuillet est collé plus bas que les précédents sur
sa feuille de montage, dont le haut est blanc. La première ligne achève
la phrase de la vue 482 (« et dont l'origine seroit à lamoitié / du côté
compris entre … ») ; ce qui est barré en tête de ligne en était une
première suite. « àlaquelle » : l'accent est écrit au-dessus de la ligne.
Dans « 40 », le 4 est repassé sur un autre chiffre. La lettre de
« (G) deviendra » est une capitale à boucle qui peut se lire G ou 6 ; les
intégrales désignées par des lettres, (C) à (I), sont celles des pages du
mémoire antérieures à ce lot. Dans la seconde valeur de $z$, $\text{sin}$ est
écrit avec son s long ; dans « $\text{cos}\frac{(\pi x m)}{A}$ » de l'avant-dernier
alinéa, le $x$ est écrit sur une autre lettre. Les trois nombres
$-m^{2}(\Pi''-2) = 10$ (vue 481), 40 et 26 s'accordent avec
$\Pi'' = 1 - 9$, $m = 1$ puis $m = 2$, et $\Pi'' = 1 - 25$, $m = 1$ : c'est
une remarque du transcripteur, non de la page.}

\page{486}

Je ne m'arrêterai pas plus longtems àla recherche des cas
renfermés dans l'intégrale (G). Mais je ne puis me dispenser
defaire remarquer encore \struck{avec quelle étonnante sagacité} que
M\textsuperscript{r} Chladni semble avoir \struck{\ill{}} \add{deviné} dans des figures composées
de lignes dues à un facteur periodique et de lignes
provenants d'un facteur different, que\struck{lles devoient
être celles qu'il étoit parmi les} \add{ces dernieres devoient être} regardées comme
equivalantes à des lignes paralleles et susceptibles de plusieurs
formes differentes. Cet apperçu est d'autant plus digne d'éloge
que ce n'est qu'après des experiences mille fois répétées,
\struck{avec} que cet habile physicien a pû apprendre à demêler,
parmi les distortions des figures, celles qui doivent être
raportées à quelqu'inégalité physique souvent inappréciable,
de celles que la théorie doit expliquer. En effet
j'ai vu bien des fois les lignes dues aux facteurs
periodiques, se séparer et se courber: mais M\textsuperscript{r} Chladni
s'est bien gardé de \struck{\ill{}} tenir compte de cet accident;
\struck{que} et malgré l'ignorance où il étoit dela théorie,
il a \add{deviné parmi les distortions qui sembloient
être de même nature, qu'elles étoient celles que l'on ne devoit
attribuer qu'à quelque \struck{inégalité} imperfection} \struck{\uncertain{senti} ne devoit pas être considerée comme
une veritable distortion dela figure qui \uncertain{la presentoit}
mais seulement comme une inégalité tenant a l'imperfection}
dela plaque vibrante. \struck{dela plaque vibrante}
\note{En tête, au milieu, « (76) », non barré ; aucun nombre en haut à
droite. C'est le dos du feuillet de la vue 485 (« (75) », « 238 ») ; le
feuillet, collé bas sur sa feuille de montage, est photographié de l'autre
côté. Devant « intégrale » une lettre est surchargée. Le mot barré sous
« deviné » (deuxième phrase) n'est pas lisible ; de même le groupe barré
après « s'est bien gardé de ». Le bas de la page est une réécriture en
interligne : les lignes d'origine (« il a senti ne devoit pas être
considerée comme / une veritable distortion dela figure qui la presentoit
/ mais seulement comme une inégalité tenant a l'imperfection / dela plaque
vibrante ») sont barrées une à une, et le texte gardé est écrit entre
elles ; on donne le texte gardé d'abord, puis les lignes barrées. « dela
plaque vibrante » est écrit deux fois, la seconde barrée. Au-dessus de
« dela figure » barré, une seconde fois « dela figure », barrée aussi. Le
bas du feuillet est blanc.}

\page{487}

A l'égard de l'intégrale (H), il est aisé de voir que quoique
les figures qui pourroient y être rapportées, doivent présenter des
lignes droites nodales dans le sens des deux axes, les sons qui
accompagneroient ces figures, devroient être proportionels
\add{a des nombres qui suivent entr'eux la même loi que ceux}
\struck{aux nombres} qui conviennent aux differens cas du mouvement
dela simple lame dont les extrémités sont appuyées
ou de celles dont les extremités seroient libres en
vertu des conditions $\frac{dz}{dx} = 0$ et $\frac{d^{3}z}{dx^{3}} = 0$.

Il est clair que les séries qui entrent dans l'intégrale
(H), ne peuvent jamais se terminer; ainsi \struck{il faudroit}
elle ne peut être d'une application facile: c'est pourquoi
\struck{que} je ne m'y arrêterai pas.

L'intégrale (I) qui appartient aux surfaces qui ne sont
périodiques que dans \struck{un} le sens des $y$, présente la
même singularité que l'intégrale (H), c'est adire
que les \struck{figur} sons qui accompagnent les figures
qui s'y rapportent, sont également proportionels
aux nombres qui conviennent aux mouvement
dela simple lame appuyée à ses deux extrémités.

La série qui entre dans cette intégrale, ne peut
non plus jamais se terminer; et il est parconséquent
impossible de déterminer les figures qu'elle peut
\note{En tête, au milieu, « (77) » barré d'un trait ; en haut à droite, à
l'encre, « 239 ». La page commence un nouvel alinéa ; elle ne continue pas
une phrase de la vue 486. La lettre des intégrales (H) et (I) est une
capitale à deux jambages, qu'on lit H entre le (G) de la vue 485 et le
(I) de cette page. L'addition « a des nombres … que ceux » est écrite dans
l'interligne au-dessus de « aux nombres », qui est barré. Le $y$ de « le
sens des $y$ » est souligné. La phrase se poursuit au-delà de cette page,
à la vue 488. Le bas du feuillet est blanc.}

\page{488}

expliquer sans avoir recours au calcul.

Cependant je crois pouvoir assurer que les fig. 79 b et
86 de M\textsuperscript{r} Chladni doivent être rapportées à cette intégrale
dans les suppositions successives. $m = 5$, $m = 6$,

En effet, si on fait $b' = 0$ \add{$n = 4$} et $m = 5$, l'intégrale (I) \struck{donnera} deviendra
\[ z = \dots \text{sin}\frac{(\pi x)}{A}\left\{a'\text{cos}\frac{(\pi x)}{A}\left\{1 - \frac{10}{3}\text{cos}^{2}\frac{(\pi x)}{A} + \frac{4}{3}\text{cos}^{4}\frac{(\pi x)}{A} + \&c.\right\} + \mathcal{A}\text{sin}\frac{(\pi 5.y)}{A}\right\}; \]
le facteur $\text{sin}\frac{(\pi x)}{A}$ donnera dabord une ligne nodale parallelle à
l'axe des $y$; et cette ligne passera par le milieu dela plaque
dont les extrémités sont libres, à ses deux limites parallelles, \add{a la ligne nodale}
on aura $\text{cos}\frac{(\pi x)}{A} = 0$, desorte que ce facteur fournira à lui
seul deux points de repos; \struck{lorsque $\text{sin}\frac{(\pi 5 y)}{A}$ sera $= 0$, et il est a}
\struck{lorsque l'on trouvera} \add{\struck{\uncertain{Donnés}} par l'intersection} des deux figures nodales dues aux suppositions successives
\struck{que j'ai déja nommé points de repos invariables}, \add{$a' = 0$ et $\mathcal{A} = 0$.} le
facteur $1 + \frac{4}{3}\text{cos}^{4}\frac{(\pi x)}{A} + \&c \ill{} \frac{10}{3}\text{cos}^{2}\frac{(\pi x)}{A}$ donnera aussi 2 points
de repos invariables, puisqu'il y a deux valeurs de $\text{cos}\frac{(\pi x)}{A}$
qui peuvent le rendre nul.

Cela posé, il devra y avoir deux flexions complettes sur
les lignes sinueuses qui composeront la figures. La fig 79 b est
en effet conforme à ces conditions: mais tandis qu'elle devroit
montrer 5 lignes sinueuses, \struck{parall} si la relation entre \uncertain{$b'$} et
$\mathcal{A}$ étoit lamême pour les deux moitiés dela plaque, il
se trouve que les trois lignes inferieures sont remplacées par
une autre figure nodales que je crois equivalantes.
\struck{\uncertain{Avant} je ne presente cette explication}
\note{En tête, au milieu, « (78) », non barré ; aucun nombre en haut à
droite : c'est le dos du feuillet de la vue 487 (« (77) », « 239 »). La
première ligne achève la phrase de la vue 487 (« les figures qu'elle peut
/ expliquer … »). « $n = 4$ » est écrit au-dessus de la ligne, entre
« $b' = 0$ » et « et ». Le coefficient $\mathcal{A}$ est la capitale bouclée
que le lot 22 a notée ainsi ; $a'$ porte son accent. Dans « libres, à ses
deux limites », le « à » est une capitale repassée sur une autre lettre,
et « libres » est repassé. L'addition « a la ligne nodale » est serrée en
bout de ligne, au bord du feuillet. Le mot qui ouvre l'interligne « … par
l'intersection » est surchargé et paraît barré. Dans le second facteur, le
signe entre « \&c » et $\frac{10}{3}$ est noyé sous une tache ; « deux »
dans « à lui seul deux points » porte au-dessus un petit astérisque sans
renvoi. « entre \uncertain{$b'$} et $\mathcal{A}$ » : la lettre porte un
accent et ressemble au $b$ de la ligne « $b' = 0$ » ; le sens voudrait
peut-être $a'$, que la page ne donne pas. « la figures », « une autre
figure nodales … equivalantes » sont écrits ainsi. La dernière ligne est
barrée. Le bas du feuillet est blanc.}

\page{489}

Si \struck{on fait ensuite} \add{continuant de} $b' = 0$, $n' = \uncertain{1}$ \struck{on et} \add{prend} $m = 6$, l'équation (I) deviendra
\[ \struck{z = \dots \text{sin}\frac{(\pi x)}{A}\left\{b'\left\{\text{cos}^{2}\frac{(\pi x)}{A} + \frac{\ill{}}{3.4.5.6}\text{cos}^{4}\frac{(\pi x)}{A} - \frac{2}{34}\right.\right.} \]
\[ z = \dots\dots \text{sin}\frac{(\pi x)}{A}\left\{a'\text{cos}\frac{(\pi x)}{A}\left\{1 - \frac{31}{2.3}\text{cos}^{2}\frac{(\pi x)}{A} + \frac{31.19}{2.3.4.5}\text{cos}^{4}\frac{(\pi x)}{A} + \frac{31.19.}{2.3.4.5.6.7}\text{cos}^{6}\frac{(\pi x)}{A}\right.\right. \]
\[ \left.\left. + \&c\right\} + \mathcal{A}\text{sin}\frac{(\pi.6y)}{A}\right\} \]
on trouvera, comme dans le cas \add{precedent}, qu'il doit y avoir une ligne
nodale \struck{qu} due au facteur $\text{sin}\frac{(\pi x)}{A}$, et que la fonction de $x$
que ce facteur multiplie, doit donner quatre points de repos
invariables, desorte que les 6 lignes sinueuses dont la figure
pourra être composée, devront également présenter 2 flexions
complettes. Mais le terme $\text{sin}\frac{(\pi 6 y)}{A}$ indique qu'il doit
y avoir 6 semblables lignes aulieu de 5 que l'on a
trouvées dans le premier cas

La fig. 86 seroit entièrement conforme à \struck{ce} résultat
si les deux courbes rentrantes qui, dans la partie superieure
de cette figure, remplacent deux des lignes sinueuses,
pouvoient être considerées comme un effet accidentel de
quelqu'inégalité delaplaque. Je serois tenté d'adopter
cette opinion, parceque, \struck{par des raisons qu'il seroit trop
long de déduire}, l'existence d'une ligne sinueuse dans
la partie de la plaque qui présente les deux courbes
rentrantes, ne permet \struck{pas} \add{guère} de supposer qu'il y ait pour
cette moitié une relation nouvelle entre $a'$ et $\mathcal{A}$.

Au reste l'application que j'ai crû pouvoir faire de
l'intégrale (I) aux deux figures 79 b et 86 n'est
\note{En tête, au milieu, « (79) », non barré ; en haut à droite, à
l'encre, « 240 ». Le feuillet est collé bas sur sa feuille de montage,
dont le haut est blanc. La page ouvre un nouvel alinéa. Première ligne :
« on fait ensuite » est barré, « continuant de » écrit au-dessus ; après
« $n' =$ » un chiffre qu'on lit 1 et qui pourrait être son 4 ; « on et »
est surchargé, et « prend » écrit au-dessus. La lettre de « l'équation
(I) » est la même capitale qu'à la vue 487. La première valeur de $z$ est
barrée d'un trait et reste inachevée ; le numérateur de sa seconde
fraction est surchargé. Dans la seconde, le dernier numérateur
« 31.19. » s'arrête sur un point. « indique » : la fin du mot est
surchargée. Le « ce » devant « résultat » est barré ; « pas » est barré
et « guère » écrit au-dessus. La phrase se poursuit au-delà de cette
page, à la vue 490 ; le bas du feuillet est blanc.}

\page{490}

que conjecturale; et il faudroit avoir recours au calcul
pour s'assurer qu'elle satisfait \add{en effet} à la description exacte
de ces figures.

En adoptant l'explication que je viens de donner,
on trouvera que le son qui accompagne la fig.
79 b, devra être proportionnel au nombre $m^{2} = 25$; et que
\struck{et celui} qui accompagnera la fig. 86, devra être proportionel
au nombre $m^{2} = 36$.

On voit que j'ai \struck{donné} \add{rendu} \struck{la description plus ou moins
detaillée} \add{compte} de plus dela moitié des figures observées par
M\textsuperscript{r} Chladni \add{dans les mouvemens des plaques quarrées.} Celles que je n'ai pas expliquées, \struck{se trouveroient} \add{pourroient être}
je crois, très aisément aussi, soit en mettant de nouvelles valeurs
dans les intégrales que j'ai données, \marginal{+ \struck{soit} en employant a la fois
ces diverses intégrales
comme je le dirai en
parlant des plaques
rectangulaires differentes
du quarré,} soit \add{enfin} à l'aide de quelqu'autre
intégrale qui ne s'est pas encore présentée, car \struck{les} figures
ne sont ni plus bizares ni plus compliquées que celles
qui ont été expliquées.

Pour complettér la comparaison entre la théorie et
l'experience, je vais m'occuper de la com\add{parai}son des sons,
\struck{qu} et quoique je n'aie encore expliqué qu'une partie
des figures, et que ce ne soit parconséquent que pour celleslà
seulement que je puisse affirmer l'équivalence entre les lignes
nodales de figures variées et les lignes nodales parallelles à
un des axes; je suivrai à cet égard l'indication de M\textsuperscript{r}
Chladni, de sorte que, lorsqu'une figure n'aura pas été
expliquée, je \struck{lui appliquerai le son} regarderai le son qu'elle
\note{En tête, au milieu, « 80 », sans parenthèses ; aucun nombre en haut
à droite : c'est le dos du feuillet de la vue 489 (« (79) », « 240 »). La
première ligne achève la phrase de la vue 489 (« n'est / que
conjecturale »). « en effet » est écrit au-dessus de « à ». Dans
« proportionnel au nombre $m^{2} = 25$ », l'exposant est écrit au-dessus
du $m$, mal séparé du mot « nombre ». « rendu » est écrit au-dessus de
« donné », « compte » au-dessus de « detaillée » ; « dans les mouvemens des
plaques quarrées » dans l'interligne au-dessus de « M\textsuperscript{r}
Chladni, Celles » ; la capitale de « Celles » est repassée. « pourroient
être » est écrit au-dessus de « se trouveroient », barré, avec deux ou
trois lettres surchargées entre les deux mots. La note de la marge
gauche est appelée par une croix après « données, » ; son premier mot,
« soit », paraît barré. Dans « la com\add{parai}son des sons », « parai »
est écrit au-dessus de lettres barrées. La phrase se poursuit à la vue
491.}

\page{491}

produit comme devant être proportionnel au \struck{un certain} nombre
que la théorie donneroit, s'il s'agissoit dela figure \struck{qui} composée
de simples lignes droites àlaquelle ce savant physicien \struck{regarde}
la croit equivalante.

N.\textsuperscript{o} 16 \emph{comparaison des sons}. si on réfléchit à l'influence que peut avoir
la moindre inégalité d'épaisseur sur les intervalles des sons qui
correspondent aux differentes figures que peuvent présenter les
plaques élastiques vibrantes, on ne pourra exiger \add{un accord rigoureux} entre les
sons indiqués par la théorie et les sons donnés par
l'experience \struck{un accord rigoureux}. Pour fixer la limite des
differences entre les résultats de la théorie et ceux de l'experience\add{,}
\add{qui devront être négligés,} \struck{qui devront être negligés} j'observerai que M\textsuperscript{r} Chladni a déterminé
lui \add{même} la quantité de l'erreur possible dans les experiences les
mieux faites: car il dit \add{\uncertain{p} 154 de son livre} « les differences que l'on pourra trouver dans
« des experiences bien faites, n'excederont jamais un demi-ton tout
au plus \struck{(V § 54.)} » Ainsi lorsqu'il s'agira de comparer entr'eux deux
sons donnés par l'experience, le rapport ne pourra jamais
differer du rapport entre les nombres auxquels la théorie
veut que ces sons soit proportionels, de plus d'un ton.
encore cette difference \struck{\ill{}} pour n'être pas hors des limites
dela probabilité, doit-elle être rare, puisqu'elle suppose
que chacune des experiences comparées, \struck{est} affectée dela
plus grande erreur possible, et que cette erreur \struck{agit} \add{influe} pour
chacune d'elles dans des sens opposés.
\note{En tête, au milieu, « (81) », non barré ; en haut à droite, à
l'encre, « 241 ». La première ligne achève la phrase de la vue 490 (« je
regarderai le son qu'elle / produit … »). « au un certain nombre » :
« un certain » est barré, et « au » peut être un « a » corrigé. Le titre
« comparaison des sons » est souligné. « un accord rigoureux », barré sur
la ligne, est récrit au-dessus de « exiger entre les » ; de même « qui
devront être négligés », barré au début d'une ligne, est récrit au-dessus
de la ligne précédente ; la virgule après « l'experience » est ajoutée à
l'endroit où la ligne est corrigée. « il dit » : le mot qui suit « il »
est surchargé, et « dit » est repassé ; au-dessus, d'une plume fine,
« \uncertain{p} 154 de son livre », où le premier signe peut être un p ou
un §. Les guillemets sont les siens, doublés en tête de la seconde ligne
de la citation ; le renvoi « (V § 54.) » qui la terminait est barré. Le
signe barré devant « pour n'être pas » est une petite capitale
surchargée. « agit » est barré, « influe » écrit au-dessus. Le bas du
feuillet est blanc.}

\page{492}

Cela posé, je vais commencer par comparer entr'eux les sons
donnés par les vibrations tournantes, c'est adire par celles
qui ne contiennent pas une periode complette, \struck{pa} mais seulement
\struck{\uncertain{deux}} demies periodes dans le sens d'un des axes.

D'après la table donnée p. 152 du traité d'acoustique
les fig. qui, \add{parmi celles que j'ai expliquées,} appartiennent à ce cas, sont les suivantes:
fig. 63, 66 a, 69--70, 74 a, 79 a--79 b. 86.

J'ai trouvé. N.\textsuperscript{o} 13 que le son de la fig. 63 doit être proportionel au nombre 2
\[ \begin{array}{llcr}
 & 66^{a} & \dots\dots\dots & 5 \\
\text{N.}^{o}\,15 \dots & 69-70 & \dots\dots\dots & 10 \\
\text{N.}^{o}\,13 \dots & 74 & \dots\dots\dots & 17 \\
\text{N.}^{o}\,15 \dots & 79^{a} & \dots\dots\dots & 26 \\
 & 79^{b} & \dots\dots\dots & 25 \\
 & 86 & \dots\dots\dots & 36
\end{array} \]
la difference entre les sons des fig. 63 et $66^{a}$ doit donc être $\frac{5}{2}$
\[ \begin{array}{lcr}
63 \text{ et } 69-70 & \dots\dots & 5 \\
63 \text{ et } 74 & \dots\dots & \frac{17}{2} \\
63 \text{ et } 79a & \dots\dots & 13 \\
63 \text{ et } 79b & \dots\dots & \frac{25}{2}. \\
63 \text{ et } 86 & \dots\dots & 18
\end{array} \]
La table P. 152 du traité d'acoustique donne
\[ \begin{array}{llcl}
\text{son de la fig.} & 63 & \dots\dots & \text{sol}.1 \\
 & 66^{a} & \dots\dots & \text{si}.2 \\
 & 69 & \dots\dots & \text{si}.3. \\
 & 70 & \dots\dots & \text{ut}^{\sharp}\,4 \\
 & 74a & \dots\dots & \text{si}^{\flat}\,4 + \text{grave} \\
 & 79a-79b & \dots\dots & \text{fa}^{\sharp}\,5 + \text{grave} \\
 & 86 & \dots\dots & \text{ut}.6.
\end{array} \]
La difference des sons qui accompagnent les fig. 63 et $66^{a}$, est donc
une octave $+$ 1 ton majeur $+$ 1 ton mineur. Cet intervalle est exprimé,
comme on sait, par $\frac{5}{2}$ c'est adire qu'il est entièrement
\note{En tête, au milieu, « (82) », non barré ; aucun nombre en haut à
droite : c'est le dos du feuillet de la vue 491 (« (81) », « 241 »). La
page commence une phrase nouvelle. Le mot barré devant « demies » est
surchargé ; on y lit « deux » avec doute. Les renvois « N.\textsuperscript{o} 13 »
et « N.\textsuperscript{o} 15 » sont à des paragraphes du mémoire
antérieurs à ce lot. Les tirets et les points qui mènent aux nombres sont
les siens, ici rendus par des points. Dans le second tableau, les
nombres de droite sont les rapports des nombres du premier au nombre 2
de la fig. 63. Dans le troisième, le chiffre qui suit la note est celui
de l'octave ; le 69 est repassé ; un petit « 6 » pâle, dans la marge
devant « $66^{a}$ », ne semble pas appartenir au tableau. La phrase se
poursuit à la vue 493 ; le bas du feuillet est blanc.}

\page{493}

conforme au résultat dela théorie.

La différence des sons qui accompagnent les fig 63 et 69, est
de deux octaves $+$ 1 ton majeur $+$ 1 ton mineur. Elle est donc
exprimée, comme le veut la théorie, par le nombre 5.

Le son qui accompagne la fig 70, devroit, d'après la théorie,
(V. N.\textsuperscript{o} \uncertain{15}), être le même que celui qui appartient à la fig. 69. Cependant
\struck{il differe} ces sons different entr'eux \add{de l'intervalle} d'un ton. J'ignore la cause
de cette difference, et je ne puis m'empêcher de remarquer combien
il est singulier que M\textsuperscript{r} Chladni ait pû se décider à
regarder \struck{ces deux fig.} comme equivalantes, deux figures
dont il sembleroit que la théorie seule peut indiquer le rapport,
puisque, par une circonstance sans doute accidentelle,
l'expérience donne, pour chacune d'elles, des sons notablement
differens.

La différence entre les sons qui accompagnent les fig. 63 et 74,
est de 3 octaves et un peu moins d'un ton et demi, c'est àdire
qu'elle est $< \frac{75}{8}$. \struck{et} En effet $\frac{75}{8}$ est $>$ que $\frac{17}{2}$ donné par la \struck{th}
théorie. \struck{et la} Ces deux rapports different entr'eux de moins
d'un ton et de plus d'un demi ton. Ainsi, \add{le son de la fig. 74 semble \uncertain{trop aigu mais}} ce cas \struck{est} renfermé \add{\uncertain{pourtant}}
dans la limite de l'erreur \struck{attribuable} \add{que l'on peut attribuer} aux experiences.
\struck{mais c'est}

La difference entre les sons qui accompagnent les figures 63 et
79 a, est de 3 octaves et un peu moins de 5 tons et $\frac{1}{2}$ ton,
c'est adire $<$ que 15; le nombre donné par la théorie est 13
qui est beaucoup trop petit: car la difference entre les deux rapports
est \struck{de plus d'un ton} entre 1 ton et 1 ton et demi, tandis
qu'il devroit être de moins d'un demi ton. Ainsi le son
\note{En tête, au milieu, « (83) », non barré ; en haut à droite, à
l'encre, « 242 ». La première ligne achève la phrase de la vue 492
(« qu'il est entièrement / conforme … »). Dans « (V. N.\textsuperscript{o} \uncertain{15}) », le 1 n'est qu'un point ou un trait très court devant le 5 ; le tableau de la vue 492 renvoie les fig. 69--70 au N.\textsuperscript{o} 15. La capitale de « Cependant »
est reprise d'un grand C en bout de ligne. « de l'intervalle » est écrit
au-dessus de « d'un ton ». Dans « $\frac{75}{8}$ », le 5 a sa queue
descendante. L'addition « le son de la fig. 74 semble \uncertain{trop aigu
mais} » est écrite très fin, sur deux étages, au-dessus de « Ainsi, ce cas
est » ; le mot « \uncertain{pourtant} » est écrit sous la ligne et entouré
d'un trait, sans signe d'insertion. « est » est barré devant
« renfermé ». « que l'on peut attribuer » est écrit au-dessus
d'« attribuable », barré. « mais c'est », au début de la ligne suivante,
est barré. La phrase se poursuit à la vue 494.}

\page{494}

qui accompagne la fig 79 a, semble être trop aigu d'environ 1 ton.
La différence est encore plus grande lorsqu'il s'agit dela fig. 79 b:
car la théorie veut que le son qui accompagne cette fig, soit
un peu plus grave que le précédent. A la verité la difference
\marginal{entre les nombres qui
appartiennent a ces
deux fig. savoir} $\frac{26}{25}$ est beaucoup moindre qu'un demi ton; et considerée
en ellemême elle pourroit être negligée; mais elle contribue
encore à eloigner la théorie de l'experience, lorsque l'on
prend le son de la fig. 63 pour terme de comparaison.

La difference entre les sons qui accompagnent les figures
63 et 86, est de 4 octaves, \struck{et} d'un ton majeur, d'un ton
mineur et d'un demi ton \struck{\ill{}} majeur. Ainsi elle est exprimée
par $\frac{64}{3}$. La théorie donne 18. \struck{ainsi} le son \struck{est plus}
de la fig 86 comparé a celui de la fig. 63. est \add{donc} plus
aigu d'un ton mineur, et $\frac{1}{2}$ demi ton majeur que ne le
veut la théorie.

Au reste on se souvient que l'explication que j'ai
crû pouvoir donner des fig. 79 b et 86, n'est qu'hypothe-
tique, desorte que les nombres qui y correspondent,
ne sont aussi que probables.

Je vais m'occuper à présent m'occuper de comparer
les sons qui appartiennent aux cas où \struck{il} les figures
sont composées d'une ou de plusieurs périodes
complettes dans le sens de chacun des deux axes.
Je commencerai par celles \add{de ces fig.} dont j'ai rendu compte
\note{En tête, au milieu, « (84) », non barré ; aucun nombre en haut à
droite : c'est le dos du feuillet de la vue 493 (« (83) », « 242 »). La
première ligne achève la phrase de la vue 493 (« Ainsi le son / qui
accompagne … »). La capitale de « A la verité » est écrite sur une autre
lettre. La note de la marge gauche, appelée par un deux-points devant
$\frac{26}{25}$, se lit à la suite de « la difference ». Devant « majeur »,
deux petits signes serrés sont barrés ; on y chercherait un dièse ou deux
f sans pouvoir le dire. « ainsi » et « est plus » sont barrés sur la même
ligne ; « donc » est écrit au-dessus de « est plus » de la ligne suivante.
Le 63 de « la fig. 63. » est repassé sur d'autres chiffres. « m'occuper »
est écrit deux fois, la première n'étant pas barrée. La phrase se
poursuit à la vue 495 ; le bas du feuillet est blanc.}

\page{495}

dans les N.\textsuperscript{os} precedens, et je prendrai pour terme de comparaison
la fig. 68 a qui est la plus simple de cette classe.

J'ai trouvé N.\textsuperscript{o} 13 que le son dela fig. $68^{a}$ doit être proportionel au nombre 8
\[ \begin{array}{llcr}
 & 71 & \dots\dots\dots & 13 \\
 & 75 & \dots\dots\dots & 18 \\
 & 80a & \dots\dots\dots & 25 \\
 & 81a & \dots\dots\dots & 29 \\
 & 82 & \dots\dots\dots & 32 \\
\text{N.}^{o}\,15 \dots & 87a,\ 87^{b}\ 88a \text{ et } 88b & \dots\dots\dots & 40 \\
\dots & 79^{a} & \dots\dots\dots & 26
\end{array} \]
la difference entre les sons des figures 68 et 71 doit donc être $\frac{13}{8}$
\[ \begin{array}{lcr}
68 \text{ et } 75 & \dots\dots & \frac{9}{4} \\
68 \text{ et } 80^{a} & \dots\dots & \frac{25}{8} \\
68 \text{ et } 81^{a} & \dots\dots & \frac{29}{8} \\
68 \text{ et } 82 & \dots\dots & 4 \\
68 \text{ et } 87-88 & \dots\dots & 5 \\
68 \text{ et } 79^{a} & \dots\dots & \frac{13}{4}
\end{array} \]
La table P. 152 du traité d'acoustique donne:
\[ \begin{array}{llcl}
\text{son de la fig} & 68 & \dots\dots & \text{si}^{\flat}\,3 - \text{aigu} \\
 & 71 & \dots\dots & \text{fa}^{\sharp}\,4 \\
 & 75 & \dots\dots & \text{ut}.5 \\
 & 80 & \dots\dots & \text{fa}^{\sharp}\,5 \\
 & 81 & \dots\dots & \text{sol}^{\sharp}.5 + \text{aigu} \\
 & 82 & \dots\dots & \text{si}^{\flat}\,5 \\
 & 87-88 & \dots\dots & \text{ut}^{\sharp}\,6 + \text{aigu} - \text{ré}\,6 - \text{aigu} \\
 & 79 & \dots\dots & \text{fa}^{\sharp}\,5 - \text{aigu}
\end{array} \]
\note{En tête, au milieu, « (85) », non barré ; en haut à droite, à
l'encre, « 243 ». La première ligne achève la phrase de la vue 494 (« dont
j'ai rendu compte / dans les N.\textsuperscript{os} precedens »). Les tirets et
les points menant aux nombres sont les siens. Dans « $\frac{25}{8}$ », le
5 est surchargé d'un autre chiffre qui le fait lire 6 ; la fraction
$\frac{29}{8}$ de la ligne « 68 et $81^{a}$ » est écrite à droite de la
colonne, faute de place. Dans « 81 » du dernier tableau, le 1 est repassé.
Le bas de la page est blanc, sous le tableau ; le raisonnement qui le
suit n'est pas sur ce côté du feuillet.}

\page{496}

La différence entre les sons qui accompagnent les figures 68 et 71,
est donc d'un peu plus de 4 tons, c'est adire $>$ que $\frac{8}{5}$ la théorie
donne $\frac{13}{8}$ qui est en effet un peu $> \frac{8}{5}$.

La différence entre les sons qui accompagnent les fig. 68 et 75,
est une octave, et un peu plus d'un ton c'est $> \frac{9}{4}$ qui est
précisément le nombre donné par la théorie

La différence entre les sons qui accompagnent les fig. 68 et
80, est d'une octave et un peu plus de 4 tons, c'est à dire
$>$ que $\frac{16}{5}$. la théorie donne $\frac{25}{8}$ qui est aucontraire $<$ que $\frac{16}{5}$:
mais ladifference entre les deux rapports, est bien audessous
dela limite des erreurs, puisqu'elle est juste un \emph{comma}.

la difference entre les sons qui accompagnent les fig. 68 et
81, est d'une octave, et \struck{beaucoup} notablement plus de 5 tons
c'est adire $>$ que $\frac{32}{9}$; et en effet le rapport donné par la
théorie savoir $\frac{29}{8}$ surpasse $\frac{32}{9}$ de plus d'un \emph{comma}.

la différence entre les sons qui accompagnent les fig. 68 et
82, est d'un peu plus de 2 octaves, c'est-a-dire $>$ que 4 qui est
le nombre donné par la théorie

la difference entre les sons qui accompagnent les fig 68 et
87--88., lesquels different entr'eux de moins d'un demiton, est
en adoptant le plus favorable àla théorie, d'une octave
et de deux tons: elle est donc exprimée par le nombre
5. donné en effet par la théorie.
\note{En tête, au milieu, « 86 », sans parenthèses ; aucun nombre en haut
à droite : c'est le dos du feuillet de la vue 495 (« (85) », « 243 »). La
page ouvre un alinéa nouveau ; elle ne continue pas de phrase de la vue
495. Dans les deux « 4 tons » (deuxième et neuvième lignes), le 4 est
écrit sur un autre chiffre. Les deux « comma » sont soulignés. Le « 68 »
des deux derniers alinéas est repassé. Un trait vertical de crayon, dans
la marge gauche, descend le long du feuillet et au-delà, sur la feuille de
montage : il n'appartient pas au texte. Le bas du feuillet est blanc.}

\page{497}

La difference entre les sons qui accompagnent les fig. 68 et
79, est une octave 3 tons majeurs et 1 ton mineur: elle est donc
exprimée par $\frac{16}{5}$, la théorie donne $\frac{13}{4}$ qui ne differe pas
sensiblement de \struck{premier} \add{ce rapport}, puisque $\frac{65}{64}$ est $<$ que le comma
$\frac{128}{125}$.

Les cas suivans sont ceux pour lesquels M\textsuperscript{r} Chladni n'a pas donné
de fig., et aussi ceux dont les fig n'ont pas été expliquées dans
les N.\textsuperscript{os} précédans. Pour les désigner, j'emploirai la notation \struck{donné}
\struck{par} imaginée par cet auteur, \struck{qui} suivant laquelle $m/n$ signifie
\struck{que}, dans \struck{le cas auquel on appliquera cette} \add{l'}expression qu'il \add{devra avoir} y \struck{aura} $\underline{m}$ lignes
droites nodales parallelles à un des axes et $n$ de pareilles lignes
parallelles à l'autre axe.

Je continuerai de comparer les sons des autres cas à celui qui accompagne
la fig $68^{a}$ -- 2/2. Cela posé, \struck{pour} la \struck{th} théorie montre que
pour .. 2/2 le son doit être proportionel a 8.
\[ \begin{array}{lcr}
4/2 & \dots\dots\dots & 20 \\
5/3 & \dots\dots\dots & 34 \\
6/3 & \dots\dots\dots & 45 \\
5/4 & \dots\dots\dots & 41 \\
6/4 & \dots\dots\dots & 52 \\
5/5 & \dots\dots\dots & 50 \\
6/5 & \dots\dots\dots & 61.
\end{array} \]
\note{En tête, au milieu, « (87) », non barré ; en haut à droite, à
l'encre, « 244 ». La page ouvre un alinéa nouveau. « premier » est barré
et « ce rapport » écrit au-dessus, avec deux ou trois lettres surchargées
devant. La notation de Chladni est écrite « m/n », les deux lettres
séparées par un trait oblique, comme « 2/2 », « 4/2 » dans le tableau.
La fin de la phrase est corrigée sur la ligne même : « dans le cas auquel
on appliquera cette » est barré d'un trait continu ; « expression »
reste, le début du mot surchargé, peut-être en « l'expression » ; « qu'il »
est suivi de « y », puis d'un mot barré (« aura » ?) ; « devra avoir » est
écrit au-dessus. On donne la lecture la plus probable sans en garantir
l'ordre. Le $m$ de « $\underline{m}$ lignes » est souligné. Les chiffres
des nombres du tableau (20, 34, 45, 41, 52, 50, 61) valent $m^{2} + n^{2}$
pour chaque couple, ce qu'on donne comme vérification du transcripteur,
non comme lecture. Le bas du feuillet est blanc.}

\page{498}

La difference entre 2/2 et 4/2 doit donc être $\frac{5}{2}$
\[ \begin{array}{lcr}
2/2 \text{ et } 5/3 & \dots\dots & \frac{17}{4} \\
2/2 \text{ et } 6/3 & \dots\dots & \frac{45}{8} \\
2/2 \text{ et } 5/4 & \dots\dots & \frac{41}{8} \\
2/2 \text{ et } 6/4 & \dots\dots & \frac{13}{2} \\
2/2 \text{ et } 5/5 & \dots\dots & \frac{25}{4} \\
2/2 \text{ et } 6/5 & \dots\dots & \frac{61}{8}.
\end{array} \]
la table p 152 donne .. son de
\[ \begin{array}{lcl}
2/2 & \dots\dots & \text{si}^{\flat}\,3 - \text{aigu}. \\
4/2 & \dots\dots & \text{ut}^{\sharp}.5 \text{ et ré}\,5 \\
5/3 & \dots\dots & \text{si}\,5 - \text{aigu et ut}\,6 - \text{aigu} \\
6/3 & \dots\dots & \text{mi}.6 \\
5/4 & \dots\dots & \text{ré}^{\sharp}.6 \\
6/4 & \dots\dots & \text{sol}.6 + \text{aigu} \\
5/5 & \dots\dots & \text{fa}^{\sharp}\,6 + \text{aigu} \\
6/5 & \dots\dots & \text{si}^{\flat} - \text{aigu}
\end{array} \]
On voit que, les deux cas qui appartiennent à 4/2 different entr'eux
d'un demiton. si on prend ré 5 -- aigu, cequi rapproche d'ut$^{\sharp}$ 5, on
trouvera que la difference entre les sons de 2/2 et de 4/2, \struck{est}
sera d'une octave et 2 tons, c'est adire qu'elle sera exprimée
comme le veut la théorie par $\frac{5}{2}$.

Pour 5/3, les deux cas different entr'eux d'un ton environ. Le plus
favorable àla théorie, savoir si 5 -- aigu, donne \struck{une diff} dans la comparaison
a 2/2, une différence de 2 octaves et 1 demi ton, Elle est donc
exprimée par $\frac{25}{6}$ ou par $\frac{64}{15}$; et $\frac{17}{4}$ qui est le rapport indiqué
par la théorie, est en effet $\frac{17}{4}$ \struck{est} $> \frac{25}{6}$ et $< \frac{64}{15}$
\note{En tête, au milieu, « (88) », non barré ; aucun nombre en haut à
droite : c'est le dos du feuillet de la vue 497 (« (87) », « 244 »). La
page commence par un tableau qui fait suite à celui de la vue 497. La
ligne « 2/2 … si$^{\flat}$ 3 -- aigu » est écrite plus petit, au-dessus de
la ligne « 4/2 », en ajout ; on la place en tête du tableau. Dans « la
table p 152 donne .. son de », la ligne entière est à gauche du tableau.
La dernière ligne de 6/5 ne donne pas le chiffre de l'octave après
« si$^{\flat}$ ». La phrase de la dernière ligne répète « $\frac{17}{4}$ »,
le second barré avec « est ». Le bas du feuillet est blanc.}

\page{499}

La difference entre les sons qui appartiennent à 2/2 et a 6/3
est d'un peu plus de deux octaves et 3 tons, c'est adire $>$ que
$\frac{50}{9}$; et en effet le rapport donné par la théorie savoir
$\frac{45}{8}$ est de fort peu $> \frac{50}{9}$.

La difference \add{entre} les sons qui appartiennent à 2/2 et a 5/4, est
d'un peu plus de 2 octaves 2 tons et $\frac{1}{2}$ \struck{demi} ton, c'est adire
$> \uncertain{\frac{128}{25}}$; \struck{Au contraire} le rapport donné par la théorie, savoir
$\frac{41}{8}$ est $<$ \struck{$\frac{128}{25}$} \add{$\frac{16}{3}$}, mais \struck{d'une quantité absolument inappréciable
a l'oreille.} \add{la} la distance entre ces deux rapports est audessous
d'un demi ton mineur

La différence entre les sons qui appartiennent a 2/2 et 6/4,
est plus de 2 octaves 4 tons et $\frac{1}{2}$ ton, c'est adire $>$ que
$\frac{20}{3}$. Aucontraire, le rapport donné par la théorie
savoir $\frac{13}{2}$ est $<$ que $\frac{20}{3}$: mais la différence entre \struck{ces} deux rapports,
est fort petite, puisqu'elle surpasse le \emph{comma} d'une quantité
inappréciable àl'oreille.

La différence entre les sons qui appartiennent à 2/2 et a
5/5, est un peuplus grande que 2 octaves et 4 tons, c'est adire
$>$ que $\frac{32}{5}$. Au contraire le rapport donné par la théorie,
savoir $\frac{25}{4}$, est $<$ que $\frac{32}{5}$: mais la difference entre ces deux
nombres est bien petite, puisqu'elle se trouve juste d'un \emph{comma}

La difference entre les sons qui appartiennent à 2/2 et à
6/5, est de 3 octaves; elle est donc exprimée par le nombre 8
qui differe de moins d'un demi ton \add{majeur} du nombre $\frac{61}{8}$ donné par la théorie.
\note{En tête, au milieu, « (89) », le 8 repassé ; en haut à droite, à
l'encre, « 245 ». La page ouvre un alinéa nouveau ; elle ne continue pas
la phrase de la vue 498. « entre » est écrit au-dessus de « les », barré ou
repassé, au deuxième alinéa. Dans « $> \uncertain{\frac{128}{25}}$ », le
numérateur est écrit sur un autre nombre (on lit 126 ou 128), et le
dénominateur 25 sur un autre ; « Au contraire » est barré, écrit lui-même
sur un autre mot. À la ligne suivante, le signe $<$ est repassé sur un
autre signe, et le même nombre « $\frac{128}{25}$ », surchargé, est
remplacé par « $\frac{16}{3}$ » écrit au-dessus. « d'une quantité absolument
inappréciable a l'oreille » est barré sur deux lignes ; « la la distance
… audessous » est écrit à droite de « a l'oreille » barré, et « d'un demi
ton mineur » dessous, en plus petit. Le « 5/5 » qui ouvre le quatrième
alinéa est tracé comme deux 5 serrés, l'oblique peu visible. « comma »
est souligné deux fois. « majeur » est écrit très fin au-dessus de la fin
de « ton », à la dernière ligne. Le bas du feuillet est blanc.}

\page{500}

Il étoit impossible d'espérer plus d'accord entre la théorie
et l'experience que n'en présente la comparaison des sons
qui accompagnent \struck{les figures} celles des figures que j'ai
expliquées, dans lesquelles il y a une ou plusieurs périodes
complettes, tant dans le sens des $\underline{x}$ que dans celui des $\underline{y}$.
Le même accord se trouve aussi dans la comparaison
des sons que, d'après l'autorité de M\textsuperscript{r} Chladni, j'ai
attribués aux cas des periodes complettes. En effet, \struck{en c'est}
en prenant \struck{a la verité} dans cette dernière classe, qui présente
d'assez grandes differences de sons rapportés aux mêmes
cas, ceux de ces cas qui s'accordent le mieux avec
la théorie, j'ai trouvé que\struck{\uncertain{lles}} les rapports indiqués
par la théorie, differoient toujours de moins d'un demiton
de ceux donnés par l'experience. On voit que cette
difference est bien audessous de celles dont les erreurs
possibles pourroient rendre compte; et que M\textsuperscript{r}
Chladni a plusieurs fois regardé comme equivalantes,
des figures qui donnoient des sons differens entr'eux
d'un demi ton, d'un ton, ou même d'une tierce
(V. latable p. 152).

Cependant il se présente ici une difficulté: Lorsqu'aulie\add{u}
de comparer entr'eux les sons donnés par les figures
qui renferment des periodes complettes dans le sens des $x$
\note{En tête, au milieu, « (90) », non barré ; aucun nombre en haut à
droite : c'est le dos du feuillet de la vue 499 (« (89) », « 245 »). La
page ouvre un alinéa nouveau. Le $x$ et le $y$ de « tant dans le sens des
$\underline{x}$ que dans celui des $\underline{y}$ » sont soulignés. Les deux
groupes barrés, « en c'est » en fin de ligne et « a la verité », ne se
lisent qu'à demi. Après « j'ai trouvé que », deux lettres barrées. La
dernière ligne s'arrête au bord du feuillet sur « des $x$ » ;
« Lorsqu'aulie\add{u} » est coupé par le bord droit. La phrase se poursuit
au-delà de ce lot, à la vue 501. Le bas du feuillet est blanc.}

\end{document}
